\chapter{Conclusão}\label{cap:conclusao}

Esta dissertação partiu de uma observação simples sobre a reconstrução do estado de jogo. De entre
os nove módulos que compõem a solução de referência do desafio SoccerNet, os dois responsáveis pela
identificação de cada jogador, a leitura do número de camisola e a atribuição da equipa e do lado do
campo, destacam-se pelo desempenho reduzido, e são precisamente os que menos atenção têm recebido
na literatura. O trabalho concentrou-se neles e alterou quatro dos nove módulos. O leitor de
números foi substituído, os módulos da equipa e do lado foram reescritos e a reidentificação ganhou
um descritor de cor, sem mudar o que já produzia. A deteção, o seguimento, a
estimativa dos pontos do campo, a calibração e a classificação de função mantiveram-se exatamente
como na solução de referência. Esta delimitação não foi apenas uma questão de âmbito, foi também uma
opção metodológica, uma vez que garante que qualquer variação de desempenho é atribuível ao que de
facto se alterou. A comparação final reforça-a. A solução de referência foi executada por nós, nos
mesmos $159$ vídeos, no mesmo ambiente e sobre os mesmos \emph{tracklets}, e o GS-HOTA foi
decomposto atributo a atributo, de modo que o ganho de cada módulo pode ser lido isoladamente. É
esta a diferença que separa este trabalho das soluções apresentadas ao desafio, que otimizam todos
os módulos em simultâneo e não permitem isolar o contributo de cada um.

Para a leitura do número, o reconhecimento ótico de caracteres da solução de referência deu lugar a
um modelo multimodal de visão e linguagem, quantizado a 4 bits para caber nos 6~GB de memória
gráfica disponíveis. A contribuição não está apenas na substituição do modelo, mas na forma como o
problema lhe é apresentado. Em vez de lhe entregar o recorte completo do jogador, isola-se
previamente a região dorsal por estimativa de pose, descartam-se os recortes sem informação
utilizável e decide-se o número de cada \emph{tracklet} apenas com as leituras em que o próprio modelo
indica estar seguro, o que se mede pela probabilidade que o modelo atribui à resposta. Foram estes três passos, e não
a escolha do modelo em si, que se revelaram determinantes. O último passo nasceu de um diagnóstico ao
próprio módulo, cuja primeira versão tratava todas as leituras como igualmente fiáveis. Reduziu os
números errados a cerca de um terço, de $12.3\%$ para $4.4\%$ das deteções de jogadores, e ao mesmo
tempo aumentou os números certos.

Para a atribuição da equipa, corrigiram-se erros no agrupamento por aparência e no modo como cada
grupo é associado a um lado do campo. O agrupamento passou a aferir a sua própria qualidade antes de
ser aceite, recorrendo a um descritor da cor do equipamento, calculado sobre a região do tronco
durante a reidentificação, quando as representações de reidentificação não separam as duas equipas.
A atribuição do lado deixou de assentar num critério único e passou a resolver-se por uma cascata
de quatro critérios de exigência decrescente, que termina num desempate pela posição do
guarda-redes.

Os ganhos, medidos sobre a totalidade dos vídeos anotados e contra a solução de referência corrida
nas mesmas condições, são claros. A exatidão da leitura do número subiu de $15.3\%$ para
$42.6\%$ sobre $159$ vídeos, com melhoria em $156$ deles e nenhuma regressão. A atribuição da equipa e do lado passou de $83.5\%$ para $92.8\%$ dos
\emph{tracklets}, e o número de vídeos em que o módulo de equipa falha, avaliado por inspeção
visual, desceu de $28$ para $7$. Na métrica oficial do desafio, o
GS-HOTA subiu $16.6$ pontos no conjunto de teste, de $31.30\%$ para $47.87\%$,
$13.1$ na validação e $13.8$ no treino, e a decomposição por atributo mostra que
esse ganho vem inteiramente dos dois atributos alterados, a equipa e o número, sendo o número o que
mais pesa. O conjunto de treino nunca foi observado durante o desenvolvimento nem contribuiu para a
afinação de qualquer limiar, e os ganhos que aí se obtêm são da mesma ordem dos do conjunto de
teste, o que afasta a objeção mais previsível a um trabalho que afina limiares sobre parte dos dados
de avaliação.

A comparação produziu ainda um resultado que não estava previsto. Ao verificar que o GS-HOTA
calculado apenas com a posição, que deveria ser igual nos dois sistemas, não o era, identificou-se
um defeito na configuração distribuída da solução de referência, que calibra a câmara em lotes de
oito \emph{frames} e aplica a homografia do primeiro aos restantes sete. O defeito custa-lhe cerca
de $5$ pontos de GS-HOTA calculado apenas com a posição e $2.4$ na métrica
completa. Todos os valores aqui reportados comparam o sistema proposto
com a referência corrigida deste defeito, o que reduz o ganho reclamado ao que é de facto devido às
alterações introduzidas.

Delimitar o alcance destes ganhos ocupou uma parte substancial do trabalho. Três tentativas de
melhorar a leitura do número, uma ponderação das leituras por qualidade, um realce de contraste e a
apresentação conjunta de todas as vistas de um \emph{tracklet}, não produziram melhorias e em alguns
casos degradaram resultados que já estavam corretos. As três chocaram com o mesmo obstáculo, o de
que nos casos difíceis o recorte entregue ao modelo já não contém o número a ler. O que funcionou foi
outra via, que passa por usar a confiança do modelo para separar as leituras em que viu o número das leituras em
que o adivinhou. Duas correções testadas depois, a fusão de fragmentos do seguimento e a imposição de
números únicos dentro de cada equipa, melhoravam indicadores intermédios mas não o GS-HOTA, e foram
por isso rejeitadas. O erro que resta deve-se sobretudo a recortes sem resolução e a trocas
de identidade do seguimento, e encontra-se nas etapas que alimentam o módulo.

No módulo de atribuição da equipa a situação é outra, e deve ser descrita tal como é. A
inspeção visual dos $159$ vídeos de cada sistema caracterizou as falhas que subsistem em três
grupos. O maior, vinte vídeos do conjunto de validação, não é deste módulo, uma vez que uma equipa joga de
amarelo e o classificador de função, que não foi alterado, dá-a como equipa de arbitragem nos dois
sistemas, retirando-a do agrupamento antes de ele começar. É este grupo que explica a distância
entre a validação e os outros dois conjuntos. O segundo grupo era, até tarde no
desenvolvimento, a única regressão própria deste trabalho, onze vídeos em que o agrupamento terminava
com quase todos os jogadores num só grupo, afirmando que todos pertenciam à mesma equipa. A causa
foi identificada nos registos de execução, a verificação de equilíbrio da partição só era aplicada a
um dos ramos da auditoria, e a correção pôde ser feita e avaliada. O módulo passa a rejeitar
qualquer partição cujo grupo menor não chegue a $8\%$ das deteções e a reagrupar pela cor do
equipamento medida nos \emph{frames}, resolvendo sete dos onze vídeos sem alterar nenhum dos
restantes $148$. Dos quatro que não se resolvem, um mantinha já a atribuição correta e conserva-a.
Dos três que subsistem, um pertence ao bloco da equipa de amarelo, outro opõe
equipamentos verde e branco, caso em que a exclusão dos píxeis de relva remove também parte das
camisolas, e o terceiro tem a anotação invertida, pelo que passa a ser penalizado justamente por
deixar de errar. O terceiro grupo são três vídeos com o lado trocado, num dos quais um critério
intermédio da cascata decide, por margem curta, antes de a posição média ou o guarda-redes serem
consultados.

Na atribuição do lado do campo produziu-se ainda uma medição que a literatura não apresenta. Sobre a
anotação de referência dos $160$ vídeos, verificou-se que o critério da posição média, adotado pela
generalidade das soluções do desafio, se inverte em $9$ casos, ou seja em $5.6\%$ dos vídeos, ao
passo que a posição do guarda-redes nunca contradiz a anotação nos $155$ vídeos em que é observável.
O alcance deste resultado deve ser enunciado com precisão. O recurso ao guarda-redes como
referência de lado não é original deste trabalho, pois já se encontra em soluções apresentadas ao
desafio. O que aqui se acrescenta é a quantificação da frequência com que a suposição alternativa
falha, que nenhuma dessas soluções apresenta, e a avaliação isolada de cada nível da cascata de
decisão. Essa medição sugere uma limitação de conceção do próprio módulo. O desempate pelo
guarda-redes ocupa o último lugar da cascata e só é acionado quando as posições médias dos dois
grupos quase coincidem, pelo que, sempre que um critério anterior decide, ainda que incorretamente,
o guarda-redes não chega a ser consultado. Num dos casos observados, as duas equipas encontram-se
na mesma metade do campo, separadas por menos de dez metros, ao passo que os respetivos
guarda-redes distam mais de setenta. A informação que resolveria o caso estava disponível e não foi
utilizada.

Por fim, a inspeção visual revelou que em cinco dos $160$ vídeos a própria anotação de referência
tem os lados invertidos, e que a métrica de equipa, ao contar apenas os \emph{tracklets}
classificados como jogadores, não regista os erros que o classificador de função esconde. Os valores
reportados foram mantidos tal como a métrica os conta, para serem comparáveis com os de outros
trabalhos, mas qualquer leitura fina desta tarefa exige as duas fontes, a métrica e a inspeção.

Quanto a trabalho futuro, identificam-se cinco linhas concretas, por ordem de prioridade.

A primeira é estimar a cor do fundo localmente, a partir dos píxeis que rodeiam cada deteção, em vez
de excluir um intervalo fixo de tom. O descritor de cor que sustenta a auditoria de partições
degeneradas assume que o verde é relva, o que falha quando uma das equipas joga de verde, único caso
em que o reagrupamento não recupera o vídeo. A alteração é local e os módulos de equipa correm em
segundos.

A segunda é elevar o guarda-redes na cascata de decisão, associando-o a um dos grupos de jogadores
em vez de o manter como desempate final. O sistema já calcula esta informação e descarta-a antes do
agrupamento.

A terceira é tratar a equipa de arbitragem dentro do agrupamento, procurando cinco grupos, as duas
equipas, os árbitros e os dois guarda-redes, como faz a solução vencedora de 2024, em vez de
depender do classificador de função para os excluir. É a via para os vinte vídeos da equipa de
amarelo, que hoje limitam o conjunto de validação nos dois sistemas.

A quarta é cortar os \emph{tracklets} nos pontos em que o seguimento troca de jogador. A inspeção
visual mostrou casos em que um \emph{tracklet} perde o jogador que seguia, reaparece noutro da mesma
equipa e leva consigo o número do primeiro. As tentativas de corrigir estes casos do lado do número
prejudicaram a métrica. A correção pertence ao seguimento, e a leitura por janela com confiança, que
já existe, permitiria detetar o ponto de troca como uma mudança do número lido com segurança.

A quinta é reduzir o custo computacional, hoje de cerca de vinte minutos por cada trinta segundos de
vídeo. A decomposição dos tempos mostra que os módulos alterados custam menos de dois segundos e que
o leitor de números custa o mesmo que o que substituiu. O custo está na reidentificação, na leitura
de todos os recortes e na estimativa dos pontos do campo em todos os \emph{frames}. Registar a confiança do modelo logo na primeira leitura, ler o número
apenas nos \emph{tracklets} que ainda não o têm, calcular a reidentificação apenas onde o seguimento
é ambíguo e, em câmara fixa, calibrar uma única vez, são os passos que tornariam o sistema
utilizável fora do conjunto de dados, e que uma avaliação em vídeos de outra proveniência deveria
acompanhar, uma vez que as sequências do SoccerNet-GSR provêm de uma única produção televisiva.

A estas acresce a direção que a análise das soluções mais competitivas aponta de forma independente
e que a decomposição por atributo confirma, uma vez que, com a posição apenas, o sistema não passa dos
$62$ pontos no conjunto de teste, e é na deteção, no seguimento e na calibração,
precisamente os módulos que este trabalho manteve intactos, que reside o salto para valores acima
dos $60\%$. Trata-se de um âmbito distinto do que aqui se delimitou.

Por último, vale a pena registar que a direção seguida coincide com a que o estado da arte veio a
adotar. A formulação da identificação como tarefa multimodal, resolvida por um modelo de visão e
linguagem, tornou-se prática corrente entre as soluções mais competitivas na edição de 2025, tendo
sido aqui alcançada de forma independente e sob restrições de recursos consideravelmente mais
severas.
